% Einheit 5 von 5 – Vergleichen, Ordnen, Runden (bank/brueche-dezimalzahlen/e5.jsonl, zusammenbau v0.1)
%% TODO Zweigzeile Teil 1: was hier gelernt wird (Satz in Schülersprache aus dem Einheitstitel) – Einheit 5
\einheitenkopf[e5]{Einheit 5 von 5 · Vergleichen, Ordnen, Runden}
\zweigzeile{ab Kl. 5, je nach Buch bis Kl. 6 · P10 oft}
\setcounter{aufgabe}{24}

%% TODO Titel: Ich-kann-Satz fehlt in der Bank (Platzhalter: Kettenname) – Nr. 25
%% TODO Anweisung: ein Satz über der ersten Teilaufgabe fehlt in der Bank – Nr. 25
\begin{aufgabe}{Welche Form?}
\begin{teile}
\teil Welche Form hat $\frac{7}{4}$? Kreuze an.\\ \kreuz{Bruch}\\ \kreuz{Dezimalzahl}\\ \kreuz{Prozent}\\ \kreuz{Potenz}
\end{teile}
\end{aufgabe}

%% TODO Titel: Ich-kann-Satz fehlt in der Bank (Platzhalter: Kettenname) – Nr. 26
%% TODO Anweisung: ein Satz über der ersten Teilaufgabe fehlt in der Bank – Nr. 26
\begin{aufgabe}{Vergleichen und Ordnen}
%% TODO \rechenplatz{n}: Schrittzahl des Grundfalls fehlt in der Bank (nur bei fester Schreibform, 2.3 b)
\begin{teile}
\teil Hänge Nullen an, bis beide gleich viele Stellen haben: $0{,}7$ und $0{,}42$. \leerfeld und \leerfeld
\teil Welche Zahl ist größer: $0{,}47$ oder $0{,}52$? \leerfeld
\teil Welche Zahl ist größer: $0{,}8$ oder $0{,}79$? \leerfeld
\teil Ordne von klein nach groß: $0{,}4$; $0{,}38$; $0{,}405$; $0{,}39$. \leerfeld $<$ \leerfeld $<$ \leerfeld $<$ \leerfeld
\teil Runde $3{,}864$ auf Zehntel. \leerfeld
\teil Welche Zahl liegt genau in der Mitte zwischen $1{,}4$ und $1{,}5$? \leerfeld
\teil Setze $<$, $=$ oder $>$ ein: $\frac{3}{5}$ __ $0{,}65$ \leerfeld
\teil Setze $<$, $=$ oder $>$ ein: $0{,}4$ __ $38\,\%$ \leerfeld
\teil Setze $<$, $=$ oder $>$ ein: $0{,}5^2$ __ $0{,}5$ \leerfeld
\teil Welcher Wert ist der kleinste: $5{,}5$; $0{,}55$; $0{,}5^2$; $55\,\%$? \hfill (P10 2023 OS) \\ \leerfeld
\end{teile}
\end{aufgabe}

%% TODO Titel: Ich-kann-Satz fehlt in der Bank (Platzhalter: Kettenname) – Nr. 27
%% TODO Anweisung: ein Satz über der ersten Teilaufgabe fehlt in der Bank – Nr. 27
\begin{aufgabe}{Zahl zwischen zwei Zahlen angeben (auch zwischen zwei Brüchen über Dezimalzahlen)}
\begin{teile}
\teil Gib eine Zahl zwischen $0{,}7$ und $0{,}8$ an. \leerfeld
\end{teile}
\end{aufgabe}

%% TODO Titel: Ich-kann-Satz fehlt in der Bank (Platzhalter: Kettenname) – Nr. 28
%% TODO Anweisung: ein Satz über der ersten Teilaufgabe fehlt in der Bank – Nr. 28
\begin{aufgabe}{gemischte Liste ordnen (Bruch, Dezimalzahl, Prozent, Potenz einer Dezimalzahl)}
\begin{teile}
\teil Ordne von klein nach groß: $\frac{1}{2}$; $0{,}45$; $48\,\%$; $0{,}7^2$. \leerfeld $<$ \leerfeld $<$ \leerfeld $<$ \leerfeld
\end{teile}
\end{aufgabe}

%% TODO Titel: Ich-kann-Satz fehlt in der Bank (Platzhalter: Kettenname) – Nr. 29
%% TODO Anweisung: ein Satz über der ersten Teilaufgabe fehlt in der Bank – Nr. 29
\begin{aufgabe}{Aussagen prüfen und die wahre ankreuzen}
\begin{teile}
\teil Welche Aussage ist wahr? Kreuze an.\\ \kreuz{$0{,}3 > 0{,}29$}\\ \kreuz{$\frac{1}{3} = 0{,}3$}\\ \kreuz{$3\,\% = 0{,}3$}
\end{teile}
\end{aufgabe}

%% TODO Titel: Ich-kann-Satz fehlt in der Bank (Platzhalter: Kettenname) – Nr. 30
%% TODO Anweisung: ein Satz über der ersten Teilaufgabe fehlt in der Bank – Nr. 30
\begin{aufgabe}{Vergleichen und Ordnen – Fehler finden, Begründen, Darstellungswechsel, Anwendung}
\begin{teile}
\teil Lea sagt: $0{,}28 > 0{,}3$, weil $28$ größer ist als $3$. Finde den Fehler.
\teil Warum ist $0{,}7$ größer als $0{,}69$? Begründe.
\teil Schreibe $0{,}65$ als gekürzten Bruch und in Prozent. \leerfeld und \leerfeld[\%]
\teil Beim 100-m-Lauf brauchen Ali $13{,}4$ s, Ben $13{,}38$ s und Can $13{,}45$ s. Wer ist am schnellsten, wer wird Dritter?
\end{teile}
\end{aufgabe}
